\section{Conclusion}
\label{sec:conclusion}
This survey studied agentic visual generation as goal-driven visual creation in which runtime evidence changes subsequent creation decisions. We defined and formalized this behavioral criterion, distinguished AVG from single-pass generation and fixed workflows, and introduced the L1--L5 levels to describe state-dependent control and cross-task improvement. We also proposed a six-component analytical framework comprising Goal Understanding, Specification, and Planning; Memory; Tool; Perception; Action; and Cross-Task Self-Improvement.

Using this framework, we compared agentic visual systems across image generation and editing, video and animation, three-dimensional and CAD creation, scientific visualization, structured documents, user interfaces, Web artifacts, and cross-domain workflows. We further reviewed training targets, trajectory supervision, policy initialization, multimodal feedback, reinforcement learning, experience reuse, skill abstraction, and transfer. The evaluation chapter organized evidence into artifact quality, goal and constraint satisfaction, trajectory and decision quality, and system and human-centered evaluation, while the frontier chapter summarized the major directions for expanding AVG scope, improving visual-agent intelligence, and establishing reliable evaluation.

Future work can extend this analysis to broader and more advanced agentic visual systems, including multimodal, embodied, multi-agent, and continually improving agents operating across longer tasks and more diverse environments.
